\section{Introduction}\label{sec:intro}
% REPORT WRITER: Skeleton only. Organize as an academic argument, not a numbered
% project diary.

\subsection{Background and Motivation}\label{subsec:intro-background}
% REPORT WRITER: Establish the IoT event-to-action reliability context and why
% an LLM-controlled decision layer needs measurable runtime controls.
% EVIDENCE SOURCES: docs/ongoing/project-overview.md; docs/ongoing/obid-scope.md.
% [REFERENCE NEEDED: IoT event-action systems; workflow automation; LLM agents]

\subsection{Problem Statement}\label{subsec:intro-problem}
% REPORT WRITER: State the bounded problem of dependable action selection and
% prevention of invalid or unapproved action release at the shared interface.
% EVIDENCE SOURCES: docs/ongoing/research-questions.md;
% docs/ongoing/collaboration-boundary.md.
% SCOPE WARNING: do not frame the problem as universal production safety.

\subsection{Aim}\label{subsec:intro-aim}
% REPORT WRITER: Map the aim directly to the stronger Obid decision/reliability
% layer and RQ1-RQ3 without claiming inherited middleware as the contribution.
% EVIDENCE SOURCES: docs/ongoing/project-overview.md; docs/ongoing/obid-scope.md.

\subsection{Research Questions}\label{subsec:intro-rqs}
% REPORT WRITER: Insert RQ1-RQ3 verbatim from docs/ongoing/research-questions.md.
% Do not paraphrase, reorder, or replace the inherited minimal-agent comparator.
% RQ1 (LOCKED): How accurately and consistently does the extended n8n-based
% agentic decision layer produce the expected IoT action across defined normal,
% malformed, and state-dependent test cases?
% RQ2 (LOCKED): How effectively do runtime structured-output validation, action
% policies, and Human-in-the-Loop approval prevent invalid or risky agent actions
% from reaching the shared IoT action interface?
% RQ3 (LOCKED): What reliability and latency differences are observed between
% the inherited minimal agent baseline and the extended Obid agentic workflow
% under the same defined IoT test cases?

\subsection{Scope and Delimitations}\label{subsec:intro-scope}
% REPORT WRITER: One temperature-to-fan scenario, one primary model, one
% single-agent workflow, one bounded-memory configuration, simulated fan.
% EVIDENCE SOURCES: docs/ongoing/obid-scope.md; docs/decisions.md;
% docs/ongoing/final-implementation-freeze.md.
% SCOPE WARNING: no multi-agent, validator-agent, hardware, device, model, or
% memory-strategy comparison; no production deployment claim.

\subsection{Contributions}\label{subsec:intro-contributions}
% REPORT WRITER: Separate OBID_CREATED cognition, runtime reliability, HITL,
% and repeated evaluation from YACOUB_INHERITED infrastructure and baselines.
% EVIDENCE SOURCES: docs/ongoing/collaboration-boundary.md;
% docs/ongoing/final-claim-evidence-map.md.
% EXPECTED TABLE: Contribution | Artifact/evidence | Provenance.

\subsection{AI Tool Use and Disclosure}\label{subsec:intro-ai-disclosure}
% REPORT WRITER: Reserve the course-required disclosure; describe assistance,
% human verification, audit gates, and limitations concisely.
% EVIDENCE SOURCES: docs/ongoing/ai-tool-use.md; docs/ongoing/codex-workflow.md.
% [CHECK METADATA: institution/course wording for required disclosure]

\subsection{Division of Work and Collaboration Boundary}\label{subsec:intro-division}
% REPORT WRITER: State one integrated system with two separately attributable
% contributions. Reuse and verification do not transfer authorship.
% EVIDENCE SOURCES: docs/ongoing/collaboration-boundary.md;
% docs/ongoing/yacoub-handoff.md; docs/decisions.md.

\subsection{Report Outline}\label{subsec:intro-outline}
% REPORT WRITER: Give a concise eight-chapter roadmap after the chapter content
% has stabilized. Do not narrate repository steps.
